% Generated draft: requires mathematical and source review.
\documentclass[11pt]{article}
\usepackage[margin=25mm]{geometry}
\usepackage{fontspec}
\setmainfont{Latin Modern Roman}
\usepackage{amsmath,amssymb,mathtools,unicode-math}
\setmathfont{Latin Modern Math}
\usepackage{enumitem,needspace,xurl,hyperref}
\setlength{\parindent}{0pt}
\setlength{\parskip}{5pt}
\setlength{\emergencystretch}{3em}
\newcommand{\N}{\mathbb{N}}
\newcommand{\Z}{\mathbb{Z}}
\newcommand{\Q}{\mathbb{Q}}
\newcommand{\R}{\mathbb{R}}
\newcommand{\C}{\mathbb{C}}
\newcommand{\F}{\mathbb{F}}
\newcommand{\A}{\mathbb{A}}
\newcommand{\PP}{\mathbb{P}}
\newcommand{\E}{\mathbb{E}}
\newcommand{\eps}{\varepsilon}
\newcommand{\dd}{\,\mathrm{d}}
\begin{document}
\section*{UnsolvedMath ID: 30004993}
\textbf{Source problem code:} catalogue entry
\subsection{Definitions and mathematical statement}
Let \(G\) be a finitely generated group, and fix a finite symmetric generating set \(S\). Write \(B_S(n)=\{g\in G:|g|_S\leq n\}\) for its word ball, where \(n\in\Z_{\geq 0}\). A \(G\)-set \(X\) is a set equipped with an action of \(G\); the action is faithful if \(\{g\in G:gx=x\text{ for every }x\in X\}=\{1\}\). For a nonempty \(G\)-set define its uniform Schreier growth by
\[
\operatorname{vol}_{G,X,S}(n)=\sup_{x\in X}\bigl|\{gx:g\in B_S(n)\}\bigr|.
\]
The source does not explicitly define its notation \(\operatorname{vol}_{G,X}\). This draft interprets it as the uniform Schreier growth above, with the generating set suppressed. It interprets \(\operatorname{vol}_{G,X}(n)\succeq n^2\) as the existence of constants \(c>0\) and \(n_0\geq 1\) such that \(\operatorname{vol}_{G,X,S}(n)\geq cn^2\) for every integer \(n\geq n_0\). This quadratic lower-bound property is independent of the chosen finite generating set. These interpretations require confirmation against the source.

The group \(G\) is torsion-free if \(g^m=1\), for \(g\in G\) and an integer \(m\geq1\), implies \(g=1\). It is solvable if its derived series, defined by \(G^{(0)}=G\) and \(G^{(j+1)}=[G^{(j)},G^{(j)}]\), reaches the trivial subgroup after finitely many steps. It is virtually abelian if it has an abelian subgroup of finite index.

\textbf{Conjectural statement.} For every finitely generated torsion-free solvable group \(G\), either \(G\) is virtually abelian, or every faithful \(G\)-set \(X\) satisfies a quadratic uniform Schreier-growth lower bound. More explicitly, if \(G\) is not virtually abelian, then for every faithful action \(G\curvearrowright X\) and every finite symmetric generating set \(S\), there exist \(c>0\) and an integer \(n_0\geq1\), allowed to depend on the action and on \(S\), such that
\[
\forall n\in\Z,\quad n\geq n_0\ \Longrightarrow\ \sup_{x\in X}|B_S(n)x|\geq cn^2.
\]
There is no restriction on the number of orbits or the cardinality of \(X\). In this non-virtually-abelian case, faithfulness forces \(X\) to be nonempty.

\textbf{Supplied qualifications.} The retained original passage refers to an unspecified ``previous result'' and reports that the conjecture holds for transitive actions and, more generally, actions with finitely many orbits. Here transitive means having exactly one orbit. The supplied dated assessment additionally reports results for metabelian groups and broad finite-rank or linear classes, without specifying the exact hypotheses of those classes. These are source-reported contextual claims, not independently verified results in this draft.
\subsection{Short English statement}
The conjecture says that every faithful action of a finitely generated torsion-free solvable group that is not virtually abelian has uniform Schreier growth at least quadratic. The bound takes the supremum over all starting points, and its constants may depend on the action. The supplied record labels the problem partially solved; the growth notation and the reconstruction from the original contextual passage require human review.
\subsection{Sources}
\begin{itemize}
\item Local UnsolvedMath registry, permanent ID \(30004993\), title: \emph{Quadratic Schreier Growth for Torsion-Free Solvable Groups}; exact source problem code: \texttt{catalogue entry}. The supplied record gives proposed year 2022 and does not name an individual proposer.
\item Supplied report citation: \emph{Groups and Dynamics: Topology, Measure, and Borel Structure} (2023), Oberwolfach Reports, pp. 107--169, DOI 10.4171/owr/2022/3. \url{https://doi.org/10.4171/owr/2022/3}. The supplied context describes Oberwolfach Reports as records of mathematical workshops and their open-problem discussions. The citation's year 2023 is retained alongside the DOI's 2022 component and the proposed year 2022.
\item Supplied literature reference and attribution: Le Boudec--Matte Bon, \emph{Growth of actions of solvable groups}. \url{https://arxiv.org/abs/2205.11924}.
\item Literature assessment dated exactly \texttt{2026-08-22}: the supplied curation labels the problem \texttt{partially\_solved} and reports the quadratic conclusion for metabelian groups, broad finite-rank or linear classes, and faithful actions with finitely many orbits, while describing the arbitrary torsion-free solvable case as open. This is supplied dated curation derived from the OWR record, not an independently verified literature assessment; its time-sensitive status and cited claims require independent verification.
\item The supplied statement status is \texttt{corrected\_verified}, with the supplied explanation that the statement was reconstructed and checked against the official Oberwolfach report. This describes the registry's verification claim; this draft has not independently visited or checked either external source.
\end{itemize}
\end{document}
